We would like to inquire whether our manuscript, ``Unifying cell-type markers for cross-study atlas building,'' would be of interest to \textit{Nature Genetics} as a Research Article.

Marker genes are how biologists talk about cell types. They are useful because they are practical, communicable, and testable. However, modern single-cell genomics faces a tension between the simplicity of this reporting scheme and the complexity of the experimental context that generated each marker. This is felt acutely in immunology, where related populations vary across tissue, disease, and activation context, making it difficult to decide when labels from different studies refer to the same biology and when they do not. This poses a central problem for atlas-building efforts that aim to reconcile cell-type annotations at scale.

Our manuscript addresses this problem by demonstrating how LLMs can help unify cell-type annotation at scale while preserving experimental context. We show that LLMs can accurately extract marker genes from papers and connect them to the underlying data from which they were derived. This lets us identify differently named immune cell types with overlapping marker profiles, as well as same-named cell types whose markers diverge across studies and contexts. We believe this will interest the broad readership of \textit{Nature Genetics} because it addresses a central barrier to single-cell atlas construction: reconciling cell-type annotations across studies without discarding the evidence that defines them.

Our main findings are:
\begin{itemize}\setlength{\itemsep}{0pt}\setlength{\parsep}{0pt}
\item We built \texttt{LLMarkers}, a benchmark of 1{,}560 unique cell type--gene pairs from seven single-cell RNA-seq studies, each linked to its source sentence or figure and supporting differential expression comparison.
\item Reported markers are enriched among highly ranked DEGs, but simple DEG cutoffs recover at best about half of the markers authors chose to report.
\item Among LLM-based approaches, text extraction best recovers paper-specific marker claims and their provenance; zero-shot generation and DEG-based selection do not.
\item Applied to 504 bioRxiv preprints, this approach yielded 12{,}501 marker associations across 2{,}720 cell-type names, revealed differently named immune cell types with overlapping markers as well as same-named cell types with divergent profiles, and is released as a searchable catalogue linking cell types and marker genes to experimental context.
\end{itemize}

We would be grateful for your assessment of whether this work would be appropriate for consideration at \textit{Nature Genetics}.
